
\documentclass[a4paper,10pt]{article}

% =========================
% PACOTES
% =========================
\usepackage[utf8]{inputenc}
\usepackage[T1]{fontenc}
\usepackage[brazil]{babel}
\usepackage{geometry}
\usepackage{parskip}
\usepackage{hyperref}
\usepackage{titlesec}
\usepackage{xcolor}
\usepackage{textcomp}
\usepackage{enumitem}

% =========================
% CONFIGURAÇÕES DO DOCUMENTO
% =========================
\geometry{
    top=0.9cm,
    bottom=0.9cm,
    left=1.0cm,
    right=1.0cm
}

\pagestyle{empty}

\definecolor{linkblue}{RGB}{52,101,164}

\hypersetup{
    pdftitle={Curriculo Nome do Candidato},
    pdfauthor={Nome do Candidato},
    colorlinks=true,
    linkcolor=linkblue,
    urlcolor=linkblue,
    citecolor=linkblue
}

\setcounter{secnumdepth}{0}

\titleformat{\section}
{\large\bfseries}
{}
{0em}
{}
[\titlerule\vspace{0.2ex}]

\titlespacing*{\section}
{0pt}
{0.45em}
{0.25em}

\setlist[itemize]{
    leftmargin=1.3em,
    itemsep=0.08em,
    topsep=0.12em,
    parsep=0em,
    partopsep=0em
}

\begin{document}

% =========================
% CABEÇALHO
% =========================

\begin{center}

{\LARGE\textbf{Nome Completo do Candidato}}

\vspace{-0.15cm}

Cidade, Estado -- País
\quad\textbullet\quad
+55 00 00000-0000

\vspace{0.01cm}

\href{mailto:email@exemplo.com}{email@exemplo.com}
\quad\textbullet\quad
\href{https://github.com/usuario}{github.com/usuario}
\quad\textbullet\quad
\href{https://www.linkedin.com/in/usuario}{linkedin.com/in/usuario}

\vspace{0.03cm}

\href{https://portfolio.exemplo.com}
{portfolio.exemplo.com}

\end{center}

\vspace{-0.30cm}

% =========================
% EXPERIÊNCIA PROFISSIONAL
% =========================

\section{Experiência Profissional}

\textbf{Empresa Exemplo}
\hfill
Remoto / Cidade, Estado

\textit{Desenvolvedor de Software}
\hfill
\textit{Jan/2024 -- Atual}

\begin{itemize}

    \item Desenvolvi e mantive aplicações backend utilizando Java e Spring Boot, implementando APIs REST, regras de negócio, autenticação e integrações entre sistemas.

    \item Modelei e implementei soluções utilizando bancos de dados relacionais e não relacionais.

    \item Implementei autenticação e autorização utilizando JWT e controle de acesso baseado em perfis.

    \item Desenvolvi integrações com APIs externas, serviços internos e sistemas de terceiros.

    \item Otimizei consultas SQL e processos de acesso a dados, reduzindo gargalos de desempenho.

    \item Desenvolvi testes automatizados para serviços, controllers e regras de negócio.

    \item Participei de cerimônias ágeis, revisões de código e planejamento técnico das entregas.

\end{itemize}

\vspace{0.01cm}

\textbf{Empresa Tecnologia XYZ}
\hfill
Cidade, Estado

\textit{Desenvolvedor Full Stack}
\hfill
\textit{Mar/2022 -- Dez/2023}

\begin{itemize}

    \item Desenvolvi aplicações web responsivas utilizando React, TypeScript, HTML e CSS.

    \item Implementei APIs REST utilizando Node.js e integração com PostgreSQL.

    \item Desenvolvi funcionalidades de cadastro, autenticação, filtros, paginação e dashboards.

    \item Realizei manutenção, correção de bugs, refatoração e melhoria de código legado.

    \item Trabalhei com Git, pull requests, code review e integração contínua.

\end{itemize}

% =========================
% PROJETOS
% =========================

\section{Projetos em Destaque}

\textbf{Plataforma de Comunicação em Tempo Real}
\hfill
\href{https://github.com/usuario/projeto}{GitHub}

\textit{TypeScript | React | Node.js | Socket.IO | WebRTC | PostgreSQL}

\begin{itemize}

    \item Desenvolvi uma plataforma de comunicação com canais de texto, chamadas de voz e presença de usuários em tempo real.

    \item Implementei comunicação em tempo real utilizando Socket.IO e chamadas de áudio utilizando WebRTC.

    \item Desenvolvi APIs REST para autenticação, gerenciamento de usuários, mensagens e canais.

    \item Implementei persistência de dados utilizando PostgreSQL e ORM.

\end{itemize}

\vspace{0.01cm}

\textbf{Sistema de Microsserviços para Gestão de Tarefas}
\hfill
\href{https://github.com/usuario/microsservicos}{GitHub}

\textit{Java 17 | Spring Boot | PostgreSQL | MongoDB | Docker}

\begin{itemize}

    \item Desenvolvi uma arquitetura baseada em microsserviços utilizando Java e Spring Boot.

    \item Implementei serviços independentes para autenticação, gerenciamento de tarefas e notificações.

    \item Desenvolvi comunicação entre serviços utilizando APIs REST e clientes HTTP declarativos.

    \item Utilizei PostgreSQL e MongoDB para diferentes necessidades de persistência.

    \item Documentei as APIs utilizando Swagger/OpenAPI e configurei os serviços para execução com Docker.

\end{itemize}

% =========================
% COMPETÊNCIAS TÉCNICAS
% =========================

\section{Competências Técnicas}

\begin{itemize}[itemsep=0.03em, topsep=0.08em]

    \item \textbf{Linguagens:}
    Java, TypeScript, JavaScript, Python e SQL.

    \item \textbf{Backend:}
    Spring Boot, Node.js, Fastify, Express, APIs REST, JWT e microsserviços.

    \item \textbf{Arquitetura \& Boas Práticas:}
    Programação Orientada a Objetos, SOLID, Clean Code, Design Patterns e Arquitetura Hexagonal.

    \item \textbf{Frontend:}
    React, Angular, HTML, CSS e Tailwind CSS.

    \item \textbf{Bancos de Dados:}
    PostgreSQL, MySQL, MongoDB, Redis, JPA/Hibernate e Prisma.

    \item \textbf{Tempo Real \& Integrações:}
    WebSocket, Socket.IO, WebRTC, Webhooks e integrações com APIs.

    \item \textbf{Testes:}
    JUnit, Mockito, Jest, Vitest, testes unitários e testes de integração.

    \item \textbf{DevOps \& Ferramentas:}
    Git, GitHub, Docker, Maven, CI/CD, GitHub Actions e Postman.

    \item \textbf{Metodologias Ágeis:}
    Scrum e Kanban.

\end{itemize}

% =========================
% FORMAÇÃO ACADÊMICA
% =========================

\section{Formação Acadêmica}

\textbf{Universidade Exemplo}
\hfill
Cidade, Estado -- Brasil

\textit{Bacharelado em Engenharia de Software}
\hfill
\textit{Conclusão prevista: Dez/2028}

% =========================
% CURSOS
% =========================

\section{Cursos Complementares}

\begin{itemize}[itemsep=0.03em, topsep=0.08em]

    \item \textbf{Java e Spring Boot -- Plataforma de Cursos}
    \hfill
    Concluído

    \item \textbf{Desenvolvimento Frontend com React -- Plataforma de Cursos}
    \hfill
    Concluído

    \item \textbf{Banco de Dados SQL e NoSQL -- Plataforma de Cursos}
    \hfill
    Concluído

    \item \textbf{Docker e DevOps -- Plataforma de Cursos}
    \hfill
    Em andamento

\end{itemize}

% =========================
% IDIOMAS
% =========================

\section{Idiomas}

\begin{itemize}[itemsep=0.03em, topsep=0.08em]

    \item \textbf{Português:} Nativo

    \item \textbf{Inglês:} Intermediário

    \item \textbf{Espanhol:} Básico

\end{itemize}

\end{document}

